\documentclass[11pt]{article}
\usepackage[margin=1in]{geometry}
\usepackage{amsmath,amssymb,amsthm}
\usepackage{graphicx}
\usepackage{booktabs}
\usepackage{longtable}
\usepackage{microtype}
\usepackage[hidelinks]{hyperref}
\usepackage{caption}

\theoremstyle{plain}
\newtheorem{proposition}{Proposition}
\theoremstyle{definition}
\newtheorem{remark}{Remark}

\captionsetup{font=small}
\setlength{\parskip}{0.4em}
\setlength{\parindent}{0pt}

\title{Coverage at a Finite Budget: Filling the Reachable Behavior Space with n = 10,000 Items}
\author{}
\date{}

\begin{document}
\maketitle
\section{Abstract}
The companion paper on infinite-horizon diversity asks how to generate an unbounded stream of items that never collapses onto its own past; its natural objective is a packing one --- keep every new item far from everything already made. This paper studies the dual problem. A fixed budget of n = 10,000 items must be spent so that the resulting set \emph{covers} as much of the generator's reachable behavior space as possible: we maximize the measure of the union of small $\varepsilon$-balls centered at the chosen embeddings, estimated by Monte Carlo against a reference pool of cheap draws from the generator itself. The coverage functional is monotone submodular, so greedy marginal-gain selection over a candidate stream inherits the classical (1 $-$ 1/e) guarantee, and its Monte-Carlo estimate concentrates at the Hoeffding rate; we verify all three facts numerically rather than taking them on faith (0 violations of diminishing returns in 5,000 randomized chain trials; greedy attains the exact optimum on all 10 enumerable instances; the empirical 95\% error of the pool estimate sits inside the Hoeffding band at every pool size tested). On two simulated domains at n = 10,000 --- red-team evaluation suites and user-persona corpora for dialogue testing --- coverage-greedy selection covers 0.545 and 0.544 of a held-out 100k reachable pool at the calibrated mid radius, against 0.508 and 0.511 for naive iid and 0.388 and 0.367 for max-min packing, while max-min wins the min-gap metric by 20--35\% and pays for it with a 3.2\% / 2.9\% junk rate against 0.0\% for the gated policies: covering and packing are different objectives and optimizing either one visibly sacrifices the other. Recursive axis refinement roughly doubles coverage of the full latent space (0.381 vs 0.214; 0.538 vs 0.272) while \emph{reducing} coverage of the fixed pre-refinement pool, which is the paper's central accounting problem: expanding the reachable space grows the denominator of the coverage ratio, so one run admits three defensible coverage numbers and we argue all of them must be reported. We also measure a \emph{reachability gap} --- the 17.4 points of coverage that no selection rule can recover because the proposal distribution cannot reach them, 92.5\% of which recursive refinement does recover --- which follows from having an embedding oracle but no inverse. A live pilot on gpt-5.6-luna (50 items) first produced a negative result -- the method covered less of a pooled reference than a naive baseline (0.231 vs 0.422) -- which we treated as a bug report rather than a finding. Diagnosis identified three measurement errors: a denominator that encoded the naive policy's own distribution rather than a space, an epsilon calibrated on a different set than it was applied to, and a comparison that conflated generation with selection. We validate the estimator against analytically known answers, then rerun budget-matched on real gpt-5.6-luna corpora with a stratum-balanced reference and a held-out scoring half (an audit of our own first result found the selector had been optimizing the evaluation set): coverage-greedy selection then beats all five literature baselines -- SemDeDup, farthest-point/k-center, Self-Instruct's ROUGE-L filter, greedy MAP-DPP, and uniform random -- by 1.043x on DALL-E instructions and 1.091x on psychometric items, down from 1.08x and 1.29x before the leak was closed. Pushed further, to released corpora (Alpaca, PersonaHub, WizardLM), the coverage metric itself breaks down: across corpora with different intrinsic scales, covered fraction at a fixed radius correlates -0.991 with within-corpus spacing and -0.986 with Vendi, ranking the most degenerate corpus first. We conclude that coverage is well-posed within a shared candidate pool and ill-posed across corpora, and report the scale-free comparison instead, where axis conditioning raises centered Vendi 16.3 -> 40.3 and eliminates duplication (19.8\% -> 0.0\%) against our own baselines while reaching about 85\% of the released corpora's Vendi at 1/100th their scale. The same data shows the covering/packing dissociation at full strength: k-center maximizes min-gap (0.696) and finishes last on coverage (0.096, a fifth of random), and the two highest-Vendi selectors are the two worst covering ones. Naive psychometric prompting is 73.6\% exact duplicates (one item repeated 2,726 times in 10,000 generations), which high temperature barely dents (71.0\%) and conditioning nearly eliminates (0.0-4.7\%) -- direct evidence that conditioning, not temperature, buys coverage, and a demonstration that nearest-neighbor epsilon-calibration is undefined on such a corpus (median self-distance exactly 0). Finally, literal and latent diversity move in opposite directions as the corpus grows (distinct-2 falls 0.944 $\to$ 0.770 while mean-centered Vendi rises 3.98 $\to$ 37.06), and the uncentered Vendi is compressed 2.7$\times$ by the embeddings' shared mean direction --- so both levels, and both centerings, need reporting.

\section{Introduction}
A team building an evaluation suite for a language model does not get an infinite horizon. They get a budget --- ten thousand prompts, say --- and a product requirement that is naturally a \emph{coverage} requirement: the suite should contain an example near every qualitatively distinct behavior the system under test might encounter, because a suite that clusters into a few families gives false assurance about everything outside them. The same shape of problem appears when generating synthetic user personas to exercise a dialogue system: near-duplicate personas waste budget, and bugs live in the parts of user-behavior space the corpus never visits.

The parent paper in \texttt{research/infinite\_horizon\_diversity/} (this work lives in its \texttt{coverage/} subdirectory and imports its modules directly) studies the open-ended version of this problem: an unbounded stream, an objective that combines typicality (mean distance to prior embeddings) with max-min novelty, and a method stack --- generator-elicited latent axes, spectral orthogonalization against the corpus second-moment matrix, recursive refinement of saturated cells, and an append-only attractor ledger --- that keeps the stream from collapsing at any horizon. This paper keeps the method stack and swaps the objective, because the budgeted problem is \emph{not} the truncation of the infinite one:

\begin{enumerate}
  \item \textbf{Max-min packing is the wrong objective at a finite budget.} Packing pushes points apart; its optima sit on the boundary of the reachable region and over-invest in extremes and outliers, because an outlier is by construction far from everything. Coverage weights regions by how much probability mass they hold and fills the bulk first. The two objectives correspond to two different classical problems --- k-center versus maximum coverage / facility location --- and their optima genuinely differ (\S{}3.4, \S{}5).
  \item \textbf{A known n permits planning.} With an unbounded horizon there is nothing to allocate; with n = 10,000 known in advance one can ask how the budget should be split across cells of the behavior space --- equally, proportionally to measure, or adaptively --- and the answer depends on the radius $\varepsilon$ at which coverage is demanded (\S{}3.5, \S{}6).
  \item \textbf{Coverage needs a denominator.} ``We covered 62\% '' is meaningless without saying 62\% \emph{of what}. We measure coverage against the generator's own reachable distribution, approximated by a pool of cheap draws --- and this choice has a sting: the sibling paper's recursive refinement \emph{expands} the reachable space, so the better the method is at opening new territory, the larger the space it is graded against (\S{}7).
\end{enumerate}

Everything is built to the same scaling discipline as the sibling paper: per-item cost is O(1) in n (fixed-size reference pool, incremental covered masks, bounded history subsamples), so nothing in the pipeline gets slower at item 10,000 than at item 100.

We make five contributions. First, a formalization of budgeted corpus construction as Monte-Carlo maximization of a union-of-balls coverage functional, with the relevant guarantees stated precisely and \emph{checked numerically} (\S{}3, \S{}6), including the no-inverse-oracle propositions and a direct measurement of the reachability gap they imply (\S{}2, \S{}6.6). Second, a simulation study on two ground-truth mixture worlds --- red-team suites and persona corpora --- comparing naive iid sampling, max-min packing, coverage greedy, gated coverage, and gated coverage with recursive refinement at n = 10,000 (\S{}5). Third, a coverage adaptation of the sibling paper's generation calculus in which the four-factor axis-scoring product collapses to a single estimable scalar --- expected marginal $\varepsilon$-ball gain of conditioning --- with an ablation isolating what calculus-guided conditioning buys (\S{}4, \S{}5.3). Fourth, a live pilot of the full method (elicited axes, calculus-guided focus, judge gate, ledger, refinement, coverage-greedy selection) on gpt-5.6-luna with local embeddings (\S{}8). Fifth, an explicit treatment of the honest-denominator problem that refinement creates (\S{}7).

\section{Problem statement}
Fix an embedding map $\varphi$ into $R^{D}$ (we use unit-normalized embeddings, D = 768 live, D = 32 in simulation) and a radius $\varepsilon$ > 0. The generator, prompted in whatever ways the pipeline can reach, induces a \emph{reachable distribution} $\mu$ over $R^{D}$: the pushforward of everything the generator can actually be made to produce. Given a budget n, choose a set S = \{$x_{1}$, \ldots{}, $x_{n}$\} of generated items to maximize the covered measure

\[
F_{\mu ,\varepsilon }(S) = \mu ( \bigcup _{x\in S} B(x, \varepsilon ) ),
\]

the probability that a fresh draw from the generator lands within $\varepsilon$ of some chosen item. Equivalently, S should be an $\varepsilon$-net for as much of $\mu$ as the budget allows. Three deliberate choices:

\begin{itemize}
  \item \textbf{Coverage of the reachable space, not of $R^{D}$.} Balls in empty space cover Lebesgue volume but no behavior. Measuring against $\mu$ makes the metric mean ``fraction of what the generator can do that the corpus has an exemplar of'' --- which is the product requirement for an eval suite. The cost of this choice is the denominator problem of \S{}7.
  \item \textbf{$\varepsilon$ is a resolution parameter, not a nuisance.} Small $\varepsilon$ asks for exemplars of fine behavioral distinctions; large $\varepsilon$ only for one exemplar per coarse region. Every result below is reported at three radii calibrated per domain (\S{}5.1), because the ranking of methods changes with $\varepsilon$.
  \item \textbf{Quality is a constraint, not part of the objective.} An item that covers new territory but fails a quality/typicality bar is not an asset; \S{}5.4 quantifies what happens when the constraint is dropped.
\end{itemize}

\textbf{The oracle asymmetry.} One assumption structures everything downstream: we have an embedding oracle E: text $\to$ $R^{D}$, and \emph{no inverse}. There is no $E^{-1}$ that decodes a target point back into text. We can compute exactly where in embedding space the next item should land --- the largest uncovered region, the trailing eigenspace --- but we cannot manufacture an item that lands there. The only handles a pipeline has are (a) language-valued latent conditioning that \emph{steers} the generator's proposal distribution, and (b) selection among sampled candidates. This is why the architecture is propose--measure--select rather than solve, and it has two consequences we state as propositions:

\begin{proposition}[no direct construction]
Coverage-optimal point sets are not directly constructible. The greedy (1 $-$ 1/e) guarantee (\S{}3.2) holds relative to the best subset of the \emph{proposal distribution's support}, not of $R^{D}$; the difference between what an unconstrained selector could cover and what a proposal-limited selector can is a \emph{reachability gap}, which is a property of the generator-plus-conditioning stack, not of the selection algorithm. We measure it directly in \S{}6.6.
\end{proposition}

\begin{proposition}[conditioning as surrogate inverse]
Language-valued conditioning is the only available surrogate for $E^{-1}$: choosing which axis to condition on next is choosing which \emph{slice} of the reachable manifold the next candidates will be sampled from, so a calculus that ranks axes by where their slices land (\S{}4) is precisely an approximate inverse --- it maps ``where we want mass'' to ``which words to condition on''.
\end{proposition}

\section{Theory}
\subsection{The coverage functional is monotone submodular}
For any measure $\mu$ and radius $\varepsilon$, define F(S) = $\mu$($\bigcup$\_\{x$\in$S\} B(x, $\varepsilon$)) on finite S $\subset$ $R^{D}$. F is monotone (adding a ball cannot uncover anything) and submodular: for S $\subseteq$ T and any x,

\[
F(S \cup  {x}) - F(S) \geq  F(T \cup  {x}) - F(T)
\]

\emph{Proof sketch.} The marginal gain of x given S is $\mu$(B(x, $\varepsilon$) \textbackslash\{\} $\bigcup$\_\{y$\in$S\} B(y, $\varepsilon$)). Since S $\subseteq$ T implies $\bigcup$\_S B(y,$\varepsilon$) $\subseteq$ $\bigcup$\_T B(y,$\varepsilon$), the set being subtracted only grows, so the gain only shrinks. This is the standard coverage-function argument; nothing about balls is essential --- any map from items to measurable ``footprints'' yields a submodular F. \qedsymbol

The same argument applies verbatim to the Monte-Carlo estimator we actually optimize. Given a reference pool P = \{$p_{1}$, \ldots{}, $p_{P}$\} of iid draws from $\mu$,

\[
\hat{F}(S) = (1/P) \cdot  \#{ p \in  P : \mathrm{dist}(p, S) \leq  \varepsilon  }
\]

is itself a finite coverage function --- item x covers the fixed subset $N_\varepsilon $(x) = \{p : $\|$p $-$ x$\|$ $\leq$ $\varepsilon$\} of the pool, and $\hat{F}$(S) = |$\bigcup$\_\{x$\in$S\} $N_\varepsilon $(x)|/P --- so $\hat{F}$ is monotone submodular \emph{exactly}, not merely in expectation. Our implementation check (\S{}6.1) found 0 violations of either property in 5,000 randomized nested-chain trials, as it must if the code is right.

\subsection{Greedy and its guarantee}
Because $\hat{F}$ is monotone submodular with $\hat{F}$($\emptyset$) = 0, the greedy algorithm that repeatedly adds the item of maximum marginal gain satisfies the Nemhauser--Wolsey--Fisher bound

\[
\hat{F}(S_greedy) \geq  (1 - 1/e) \cdot  max_{|S|\leq n} \hat{F}(S) \approx  0.632 \cdot  OPT,
\]

and this is tight for the class (maximizing coverage is NP-hard, and beating 1 $-$ 1/e is hard under standard assumptions, Feige 1998). \S{}6.3 checks the bound against \emph{exact} optima on instances small enough to enumerate; greedy attained the optimum itself on all ten instances, comfortably above the bound --- typical behavior, since the 1 $-$ 1/e worst case requires adversarial structure.

\subsection{Sequential generation is streaming greedy}
A generation pipeline cannot pick from a precomputed ground set: candidates arrive a few at a time (K per step) and must be accepted or discarded immediately. This is the streaming submodular maximization regime. Sieve-streaming (Badanidiyuru, Mirzasoleiman, Karbasi, Krause 2014) maintains geometric thresholds and accepts any arrival whose marginal gain clears one, achieving a (1/2 $-$ $\delta$) guarantee in one pass with O((k log k)/$\delta$) memory, independent of stream length. Our per-step rule --- take the best of K fresh candidates by marginal gain --- is the natural best-of-K relaxation: weaker in the worst case (an adversarial stream can starve it) but stronger on exchangeable streams like ours, where each batch is iid from the current proposal distribution and thresholding would only discard budget. Empirically the K = 4 stream rule retains most of full greedy's value over an identical ground set (0.484 vs 0.528 covered fraction; \S{}6.2) at a per-item cost that never touches the whole ground set. What matters for scaling is that the marginal-gain oracle is O(K $\cdot$ P) per step against a fixed-size pool with an incrementally maintained covered mask --- O(1) in n.

\subsection{Coverage is not packing}
Max-min selection (choose the candidate farthest from the accepted set) is the greedy algorithm for the k-center / packing objective, and the sibling paper shows it is the right \emph{shape} of objective for an unbounded stream once a typicality anchor keeps it on-manifold. At a finite budget against a fixed $\mu$ the two objectives pull apart:

\begin{itemize}
  \item k-center cares about the worst-covered \emph{point of S itself}; its optima push to the boundary of the support and place points in regions of vanishing measure (an outlier is far from everything, so packing \emph{rewards} it --- the novelty--junk conflation of the sibling paper, now with a budget price attached: every junk item is a ball spent covering $\approx$ 0 measure).
  \item Maximum coverage / facility-location-type objectives weight territory by measure: a second exemplar in a heavy mode can be worth more than a first exemplar of a negligible one, and junk is worth nothing because $\mu$ puts almost nothing near it.
\end{itemize}

\S{}6.2 isolates the contrast with no generator in the loop: on one shared candidate set, greedy-coverage covered 0.528 of a held-out pool where farthest-point packing covered 0.143 --- while packing's min-gap (1.44) beat greedy's (0.577) by 2.5$\times$. Neither is ``better''; they optimize different functionals, and \S{}5 shows the same double dissociation with the full generation loop in place.

\subsection{Estimation error of the Monte-Carlo coverage}
For fixed S, each pool point contributes an iid Bernoulli(F(S)) indicator, so Hoeffding/Chernoff gives

\[
P( |\hat{F}(S) - F(S)| > t ) \leq  2 exp(-2 P t^{2}),
\]

i.e. a 95\% band of $t_{95}$ = sqrt(ln(2/0.05)/(2P)) --- about $\pm$0.030 at P = 2,000 and $\pm$0.015 at P = 8,000. Two caveats we take seriously. First, the bound is for fixed S; a set \emph{selected} by optimizing $\hat{F}$ on the same pool is biased upward on that pool (the winner's curse), which is why every simulation number in \S{}5 is reported on held-out pools the selector never saw, and the live pilot reports selection-pool coverage explicitly labeled as such. Second, uniform-over-S guarantees would need a union bound over the (exponentially many) candidate sets; we do not rely on one --- held-out evaluation makes the fixed-S bound the relevant one. \S{}6.4 confirms the band empirically: at every pool size tested, $\geq$ 99\% of 200 independent pool estimates fell within $t_{95}$ of a 200k-pool ground truth.

\subsection{Allocation under a known budget}
Suppose the space decomposes into cells (modes) with measures $w_{m}$, and within-cell placement is near-saturating: with $c_{m}$ items in cell m, the covered measure of that cell is approximately $w_{m}$ $\cdot$ $g_{m}$($c_{m}$) with $g_{m}$ increasing and concave (each additional exemplar covers a shrinking uncovered remainder --- submodularity again, now per cell). Maximizing $\Sigma$ $w_{m}$ $g_{m}$($c_{m}$) subject to $\Sigma$ $c_{m}$ = n is a concave allocation problem whose optimum is the water-filling rule: equalize marginal gains $w_{m}$ $g_{m}$$'$($c_{m}$) across cells. Three closed-form rules bracket the regimes:

\begin{itemize}
  \item \textbf{$\varepsilon$ large relative to cell diameter} ($g_{m}$ saturates after \textasciitilde{}1 item): any rule that puts at least one item per cell is near-optimal; equal allocation wastes least; proportional over-invests in heavy cells.
  \item \textbf{$\varepsilon$ small} ($g_{m}$ near-linear over feasible $c_{m}$, slope $\propto$ per-item covered mass $\approx$ density $\times$ ball volume): marginal gain is $w_{m}-weighted$ throughout, and proportional-to-measure allocation dominates equal.
  \item \textbf{In between}, neither closed form matches water-filling; greedy marginal coverage \emph{is} water-filling implemented empirically, and wins at the radius it optimizes.
\end{itemize}

The measured allocation experiment (\S{}6.5) shows exactly this pattern, plus a failure mode the clean theory hides: once greedy has saturated the pool at its selection radius, its marginal signal is identically zero and the tie-breaking rule silently decides where the rest of the budget goes.

\textbf{Depth per condition and switch cost.} The sibling paper's slice theory (\texttt{../verify\_slices.py}, all checks passing) sharpens the within-cell side of allocation. For a \emph{fixed} prompt/spec the generator concentrates on an m-dimensional slice with m $\ll$ d, so fixed-prompt novelty decays as $n^{-1/m}$ (fitted slopes $-$0.472, $-$0.309, $-$0.195 for m = 2, 3, 5 --- governed by m, not by the ambient dimension), and one slice $\varepsilon$-covers a vanishing $\varepsilon ^{d-m}$ fraction of the manifold. The budgeted consequence is a breadth-vs-depth rule whose correct form (their P3, corrected from its first statement) is: with \emph{free} condition-switching the optimum is the boundary n\textbackslash\{\}\emph{ = 1 --- one sample per spec, then move --- and an interior optimal depth exists only once switching carries a real cost c (eliciting a spec, embedding levels, an occasional refinement call): their measured n\textbackslash\{\}} rises from 1 (c = 0) to 10 (c = 3) to 30 (c = 30). Optimal depth is set by the switch-cost-to-sample-cost ratio, not by $\varepsilon$ alone. Our pipelines sit near the cheap-switching regime --- in simulation a spec switch is free, and in the live pilot we measure c $\approx$ 0.10 generation-equivalents (\S{}8: 5 axis-elicitation, refinement and ledger-mining calls amortized over 50 steps, against 3 generations per step) --- which is why every policy here draws its K candidates from K \emph{fresh} specs rather than sampling any spec deeply. At that switch cost the sibling's sweep puts the optimum at or adjacent to n\textbackslash\{\}* = 1, which is what we do.

\section{Method}
The pipeline is the sibling paper's stack with the selection objective and its bookkeeping swapped from packing to covering. Concretely:

\textbf{Reachable pool (the denominator).} Before selection begins, draw a pool of cheap, unconditioned samples from the generator and embed them. This pool \emph{is} the Monte-Carlo measure: coverage means coverage of it. In simulation the pool has 20,000 points (selection) plus held-out pools of 20,000 (tracking) and 100,000 (final evaluation); live, 240 one-shot prompts. Cost is O(P) once, amortized over the whole run.

\textbf{Elicited axes and specs.} The generator is asked --- before generating any items --- to name the latent axes along which items in this domain can differ, each with discrete language-valued levels (live: 6 axes for evaluation prompts, e.g. persona pressure, trap construction, response obligation). A spec is one level per applicable axis; K specs are sampled per step and one candidate is generated under each. Axes give the proposal distribution spread the base generator lacks; they are the reason the method's candidates are worth ranking at all.

\textbf{Calculus-guided conditioning (the coverage adaptation).} The sibling paper's generation calculus (\texttt{../calculus.py}) ranks candidate axes by a four-factor product --- spread (are the axis's levels different from each other), transversality (does their variation point outside the corpus's occupied eigenspace), independence (measured against axes already in use, never against rival candidates --- scoring candidates against each other makes duplicates zero each other out), and headroom (an entropy deficit over the axis's sampled levels). Under a \emph{coverage} objective this product collapses, and the collapse is itself a packing-vs-covering result. Packing has no per-axis marginal decomposition --- min-gap is a global property --- so the calculus must approximate ``will this axis move the slice somewhere new'' from four separate geometric proxies. Submodular coverage has an exact marginal quantity: the \textbf{expected marginal $\varepsilon$-ball gain of conditioning on the axis}, $E_{x\sim axis}$[F(S $\cup$ \{x\}) $-$ F(S)]. That one scalar subsumes the four factors: an axis with near-synonym levels (no spread) or whose slices sit inside covered territory (no transversality) or which duplicates an axis already exploited (no independence, because the covered mask already contains that axis's contribution and re-scores every candidate after every accept --- the same greedy re-scoring insight as the calculus's \texttt{select\_axis\_set}, executed by the objective itself) or whose region is already covered (no headroom) all have small expected marginal gain, and the headroom that remains is \emph{measure-weighted} by construction --- a large under-sampled region beats a small one because it holds more uncovered pool mass. In simulation we estimate the scalar directly: each spec carries an estimate initialized from an 8-draw probe at creation and updated as an EMA of observed candidate gains, and the K specs per step are drawn proportionally to it with a 20\% uniform exploration floor (\texttt{gated\_coverage\_refine\_calc}). In the live pilot, where axes rather than opaque specs are the unit, we keep the sibling's embedding-based spread and transversality (level-description embeddings against the accepted corpus's occupied eigenspace) and replace its entropy headroom with the measure-weighted estimate (per-level mean observed marginal gain, optimistic for unused levels); the top-two axes by this promise become the step's dominant design pressures, with their most-different levels chosen by farthest-point selection in level-embedding space --- the calculus acting as the surrogate inverse of Proposition 2.

\textbf{Coverage-greedy selection.} Each step: embed the K candidates, compute each one's marginal gain --- the number of not-yet-covered pool points within $\varepsilon _sel$ --- and accept the gated candidate of maximum gain (ties broken by judge score). The covered mask is updated incrementally; per-step cost is O(K $\cdot$ P + P) regardless of n.

\textbf{Judge gate (quality/typicality).} A separate LLM call (generation never self-grades) scores each candidate on a rubric; candidates below the bar are ineligible regardless of gain. In simulation the gate also enforces a typicality z-score against the running centroid, mirroring the sibling paper's anchor. The gate is negative supervision only: if every candidate fails, the step falls back to best-judged rather than emitting nothing, and the rejection is logged.

\textbf{Attractor ledger.} Every 20 accepted items (live; 50 in the sibling paper), the accepted corpus is shown to the model with the question ``what do these have in common?''; the mined attractors are appended to a JSONL ledger, and a bounded top-K slice becomes explicit negative constraints in subsequent generation prompts. The ledger is append-only; its influence on any single prompt is fixed-size.

\textbf{Recursive refinement + pool refresh.} When recent accepted items' marginal gains saturate (the region's balls mostly re-cover covered ground), the responsible region of spec-space is described back to the generator, which either splits an axis level into finer sub-levels or mints a new axis meaningful inside that cell --- conditional axes forming a tree, exactly as in the sibling paper. The budgeted setting adds one obligation the infinite-horizon setting does not have: refinement changes the reachable distribution, so the reference pool must be refreshed with draws from the newly opened region, or the selector will see zero gain precisely where the new territory is (its balls cover no \emph{old} pool points). This pool-refresh step is the operational face of the honest-denominator problem (\S{}7).

\section{Simulation study}
\subsection{Worlds}
Both domains instantiate the sibling paper's MixtureWorld pattern with a spec layer added between policy and world, so that policies only ever touch the handles a real pipeline has. A world has M latent modes (mixture of anisotropic Gaussians, centers at norm 1.5, per-mode $\sigma$ $\in$ [0.10, 0.22]) with skewed Dirichlet(0.6) weights, latent per-mode quality in [0.5, 0.9], a noisy judge ($\sigma$ = 0.10) over it, and a rare diffuse junk component (weight 0.8\%, $\sigma$ = 3.0, quality $\approx$ 0.05) representing off-manifold text. The policy addresses the generator through \emph{specs}: opaque handles that map (invisibly) to modes. Crucially, base specs reach only a subset of modes --- 25 of 60 (redteam), 20 of 40 (persona) --- because a real generator's unconditioned repertoire does not span its latent space; refinement is the only way to grow reach.

\begin{itemize}
  \item \textbf{redteam} (evaluation-suite world): M = 60 attack families, full-dimensional scatter in D = 32.
  \item \textbf{persona} (dialogue-testing world): M = 40 archetypes whose centers lie on a random 8-dimensional linear patch of $R^32$ (low intrinsic dimension, as embedded persona text exhibits), with small ambient jitter.
\end{itemize}

Radii are self-calibrated per domain: $\delta$ = median nearest-neighbor distance from held-out reachable draws to n = 10,000 iid reachable draws, and results are reported at $\varepsilon$ $\in$ \{0.6$\delta$, $\delta$, 1.6$\delta$\}. By construction naive iid covers $\approx$ 50\% of the pool at $\varepsilon$ = $\delta$, which makes the mid radius interpretable on sight. Measured: $\delta$ = 0.887 (redteam), 0.787 (persona).

\subsection{Policies}
All policies run at n = 10,000 with K = 4 candidates per step (naive takes the first draw), identical world geometry, and disjoint RNG streams:

\begin{table}[htbp]
\centering
\small
\begin{tabular}{lrrrr}
\toprule
\textbf{policy} & \textbf{selection rule} & \textbf{gate} & \textbf{refinement} & \textbf{conditioning} \\
\midrule
naive & first draw & -- & -- & uniform specs \\
maxmin & argmax min-dist to accepted set (bounded 1,000-subsample) & -- & -- & uniform \\
coverage & argmax marginal $\varepsilon$-ball gain on 20k estimation pool & -- & -- & uniform \\
gated\_coverage & same & judge $\geq$ 0.35 $\wedge$ anchor z $\leq$ 3 & -- & uniform \\
gated\_coverage\_refine & same & same & spec splits + new specs + pool refresh & uniform \\
gated\_coverage\_refine\_calc & same & same & same & promise-weighted (\S{}4) \\
\bottomrule
\end{tabular}
\end{table}

\subsection{Results}
Covered fraction of the held-out 100k base-reachable pool at the calibrated mid radius (redteam $\varepsilon$ = 0.887 / persona $\varepsilon$ = 0.787), with the coverage of the full latent space (the oracle pool base policies cannot reach), junk, quality, and the packing metric, all at n = 10,000:

\textbf{redteam} ($\delta$ = 0.887):

\begin{table}[htbp]
\centering
\small
\resizebox{\linewidth}{!}{%
\begin{tabular}{lrrrrrr}
\toprule
\textbf{policy} & \textbf{covered base @$\varepsilon _mid$} & \textbf{covered world @$\varepsilon _mid$} & \textbf{junk \%} & \textbf{mean quality} & \textbf{min-gap} & \textbf{Vendi} \\
\midrule
naive & 0.508 & 0.195 & 0.70 & 0.684 & 0.430 & 23.8 \\
maxmin & 0.388 & 0.159 & 3.19 & 0.690 & \textbf{0.549} & 23.8 \\
coverage & \textbf{0.545} & 0.214 & 0.70 & 0.681 & 0.458 & 23.4 \\
gated\_coverage & 0.541 & 0.214 & \textbf{0.00} & 0.761 & 0.403 & 22.2 \\
gated\_coverage\_refine & 0.492 & 0.381 & \textbf{0.00} & 0.765 & 0.056 & 18.8 \\
\textbf{gated\_coverage\_refine\_calc} & \textbf{0.567} & \textbf{0.427} & \textbf{0.00} & 0.760 & 0.138 & 23.9 \\
\bottomrule
\end{tabular}
}
\end{table}

\textbf{persona} ($\delta$ = 0.787):

\begin{table}[htbp]
\centering
\small
\resizebox{\linewidth}{!}{%
\begin{tabular}{lrrrrrr}
\toprule
\textbf{policy} & \textbf{covered base @$\varepsilon _mid$} & \textbf{covered world @$\varepsilon _mid$} & \textbf{junk \%} & \textbf{mean quality} & \textbf{min-gap} & \textbf{Vendi} \\
\midrule
naive & 0.511 & 0.257 & 0.90 & 0.707 & 0.431 & 14.5 \\
maxmin & 0.367 & 0.191 & 2.91 & 0.705 & \textbf{0.549} & 16.3 \\
coverage & \textbf{0.544} & 0.276 & 0.51 & 0.702 & 0.443 & 14.0 \\
gated\_coverage & 0.537 & 0.272 & \textbf{0.00} & 0.786 & 0.410 & 13.1 \\
gated\_coverage\_refine & 0.501 & 0.538 & \textbf{0.00} & 0.787 & 0.064 & 9.0 \\
\textbf{gated\_coverage\_refine\_calc} & \textbf{0.571} & \textbf{0.582} & \textbf{0.00} & 0.757 & 0.102 & 11.5 \\
\bottomrule
\end{tabular}
}
\end{table}

Four observations, consistent across both domains.

\begin{enumerate}
  \item \textbf{The double dissociation is real with the full generation loop in place.} Max-min wins its own objective decisively (min-gap 0.549 vs $\approx$ 0.43--0.46 for everything else) and pays 12--18 points of covered fraction for it --- \emph{below naive iid} on the coverage objective in both domains, because its budget migrates to boundaries and outliers where the reachable measure is thin. Coverage-greedy beats naive by 3.3--3.7 points at $\varepsilon _mid$ on the held-out pool; the gap versus maxmin is 15.7 (redteam) and 17.6 (persona) points.
  \item \textbf{Radii change the ranking.} At the large radius (1.6$\delta$) every policy is $\approx$ 0.98--0.99 on the base pool --- the base reachable set is easy to cover coarsely, and differences live at and below $\delta$. At the small radius (0.6$\delta$), refinement dominates in redteam (0.031 vs 0.014 naive: tight child specs supply fine-grained exemplars) while in persona the low-intrinsic-dimension geometry makes small balls so poor at catching measure (covered fractions of $10^{-3}$) that no policy separates meaningfully --- resolution demands must respect the manifold dimension (see also \S{}6.7).
  \item \textbf{Refinement trades base-pool coverage for reach.} The refine policies give back \textasciitilde{}4 points on the fixed base pool (0.49--0.50 vs 0.54 gated) and in exchange nearly double coverage of the full latent space (0.381 vs 0.214 redteam; 0.538 vs 0.272 persona), reaching all 60/60 and 40/40 modes versus the 25/60 and 20/40 the base specs can address. Which of those numbers is ``the'' coverage of the run is exactly the denominator question of \S{}7.
  \item \textbf{Diversity metrics disagree, instructively.} The refine policies have the \emph{lowest} Vendi scores (18.8 and 9.0) while covering the most latent territory --- linear-kernel Vendi measures global spectral spread, and the dense local clusters that tight child specs produce (min-gap 0.056) concentrate the spectrum even as they open new regions. A single diversity number, whether min-gap or Vendi, cannot summarize a coverage run; the covered-fraction-by-denominator table can. This is the coverage-side face of the sibling paper's metric trap (\S{}11): its fixed-prompt policy scored \emph{best} on residual-headroom precisely because it covered nothing.
  \item \textbf{Calculus-guided conditioning dominates, and repairs refinement's regression.} Adding promise-weighted spec sampling (\S{}4) to the refine policy is the only change that wins on \emph{both} denominators at once: it is the best policy on the fixed base pool (0.567 / 0.571, ahead of plain coverage-greedy at 0.545 / 0.544) \emph{and} the best on the full latent space (0.427 / 0.582, ahead of unguided refinement at 0.381 / 0.538). It does this while recovering most of what unguided refinement gave up elsewhere: Vendi rises from 18.8 to 23.9 (redteam) and 9.0 to 11.5 (persona), and min-gap from 0.056 to 0.138 and 0.064 to 0.102. The mechanism is straightforward --- unguided refinement spends budget uniformly across child specs, most of which land in already-covered territory, whereas promise-weighted sampling concentrates draws on the specs whose expected marginal $\varepsilon$-ball gain is still large. The ablation isolates this: the two policies differ \emph{only} in how specs are sampled, with identical gates, identical refinement triggers, and identical pool refresh.
\end{enumerate}

\begin{figure}[htbp]
\centering
\includegraphics[width=\linewidth]{figures/fig_coverage_growth.png}
\caption{Figure 1. Coverage growth at n = 10,000 on a held-out 20k reachable pool, three radii per domain. Coverage-greedy dominates at the selection radius; refinement wins at small radii (redteam) and on the full latent space; max-min trails everything.}
\end{figure}

\begin{figure}[htbp]
\centering
\includegraphics[width=\linewidth]{figures/fig_tradeoff.png}
\caption{Figure 2. The covering/packing double dissociation: covered fraction at the mid radius versus min-gap, with junk rates annotated. No policy wins both objectives.}
\end{figure}

\subsection{Junk, quality, and the gate}
The junk column isolates the novelty--junk conflation at a finite budget. Max-min accepts 3.2\% / 2.9\% junk --- with K = 4 candidates and junk weight 0.8\%, essentially every junk draw that appears among the candidates is selected, because a diffuse off-manifold point maximizes distance to everything; each such acceptance is a ball spent covering $\approx$ 0 reachable measure. Ungated coverage does \emph{not} share the attraction: junk lands at or below the naive base rate (0.70\% / 0.51\%), because a ball around junk covers almost no pool mass, so junk wins the argmax only when every candidate's marginal gain is $\approx$ 0 (late-run ties). The gate removes junk entirely (0.00\% over 10,000 accepts in all four gated runs) and raises mean true quality by 0.08 in both domains at a cost of $\leq$ 0.7 points of covered fraction --- the cheap insurance the sibling paper's typicality gate promised, now priced under a budget. The residual case for the gate under coverage is therefore not junk \emph{attraction} (packing's failure) but junk \emph{indifference}: a quality-blind coverage selector will still spend late-run ties on garbage, and mean quality drifts with whatever the proposal distribution emits.

\begin{figure}[htbp]
\centering
\includegraphics[width=\linewidth]{figures/fig_spectrum_vendi.png}
\caption{Figure 3. Left, center: Vendi trajectories -- refinement's dense child clusters lower spectral diversity even as latent coverage doubles. Right: marginal gain of accepted items decays (diminishing returns in the wild); refinement events (shaded) repeatedly reset the decay for the refine policies.}
\end{figure}

\section{Numerical verification of the theory}
All checks in \texttt{verify\_theory.py}; all numbers in \texttt{figures/summary\_theory.json}.

\textbf{6.1 Submodularity (exact).} 200 candidates, 5,000-point pool, $\varepsilon$ = 1.0. Over 5,000 random nested chains S $\subset$ T with a fresh x: 0 violations of diminishing returns, 0 of monotonicity, worst violation 0.0.

\textbf{6.2 Selection rules head-to-head (no generator).} One shared ground set of 2,000 candidates, budget k = 300, $\varepsilon$ = 1.0, held-out 20k pool:

\begin{table}[htbp]
\centering
\small
\begin{tabular}{lrr}
\toprule
\textbf{rule} & \textbf{covered fraction (held-out)} & \textbf{min-gap} \\
\midrule
full greedy & 0.528 & 0.58 \\
stream greedy, K = 4 & 0.484 & 0.68 \\
random & 0.420 & 0.53 \\
max-min (Gonzalez) & 0.143 & 1.44 \\
\bottomrule
\end{tabular}
\end{table}

Packing is catastrophic \emph{as a coverage algorithm} (below random by 3$\times$) while being unbeatable at its own metric --- the pure-form double dissociation. Stream greedy with K = 4 keeps 92\% of full greedy's coverage.

\textbf{6.3 The (1 $-$ 1/e) bound against exact optima.} Ten instances with |ground| = 30, k = 5 (142,506 subsets enumerated each): greedy/OPT ratio was 1.0 on every instance; bound 0.632 never approached.

\textbf{6.4 Monte-Carlo error.} Fixed S of 300 items, ground truth from a 200,000-point pool (F = 0.4134). Over 200 independent pools per size: empirical 95th-percentile absolute error 0.042 / 0.020 / 0.011 at P = 500 / 2,000 / 8,000, versus Hoeffding $t_{95}$ = 0.061 / 0.030 / 0.015; the fraction of estimates inside the band was 0.99--0.995 $\geq$ 0.95 everywhere. P = 20,000 (our estimation pool) implies $t_{95}$ $\approx$ 0.0096: pool noise is well below every effect size we interpret.

\textbf{6.5 Allocation rules.} n = 10,000 across the redteam world's M = 60 cells with \emph{known} measures (the oracle planning question), evaluated against a 50k full-mixture pool:

\begin{table}[htbp]
\centering
\small
\begin{tabular}{lrrrr}
\toprule
\textbf{rule} & \textbf{$\varepsilon$ = 0.6} & \textbf{$\varepsilon$ = 1.0} & \textbf{$\varepsilon$ = 1.6} & \textbf{$\varepsilon$ = 2.4} \\
\midrule
equal per cell & 0.046 & 0.494 & 0.9916 & 0.9921 \\
$\propto$ measure & \textbf{0.060} & 0.547 & 0.9917 & 0.9921 \\
$\propto$ sqrt(measure) & 0.055 & 0.531 & 0.9918 & 0.9921 \\
greedy marginal (at $\varepsilon$ = 1.0) & 0.004 & \textbf{0.618} & 0.9919 & 0.9921 \\
\bottomrule
\end{tabular}
\end{table}

Three regimes, as \S{}3.6 predicts. At large $\varepsilon$ every rule saturates (all $\approx$ 0.992 --- the remaining 0.8\% is junk measure no allocation can reach). At the greedy-optimized radius, greedy beats the best closed form by 7 points: water-filling in the wild. At small $\varepsilon$, proportional wins among closed forms (near-linear per-cell returns favor measure-weighting) --- and greedy at $\varepsilon$ = 1.0 is \emph{terrible} (0.004), for an instructive reason: greedy saturated its $\varepsilon$ = 1.0 pool after roughly half the budget, its marginal signal went identically to zero, and argmax tie-breaking silently dumped 5,083 of 10,000 items into one arbitrary cell (correlation of greedy's allocation with every closed form $\leq$ 0.17). Lesson for practitioners: greedy's guarantee says nothing about what it does \emph{after} it has won; at a finite budget you must either select at the smallest $\varepsilon$ you care about, or hand the post-saturation residual to an explicit rule (e.g. proportional), or lower $\varepsilon$ adaptively as gains vanish.

\begin{figure}[htbp]
\centering
\includegraphics[width=\linewidth]{figures/fig_allocation.png}
\caption{Figure 4. Left: allocation rules at four radii -- proportional wins at small radius, greedy at its own radius, everything saturates at large radius. Right: greedy's realized allocation is measure-aware but saturating, with the post-saturation tie-breaking pathology visible as a single 5,083-item spike.}
\end{figure}

\textbf{6.6 The reachability gap (Proposition 1, measured).} Identical greedy selection ($\varepsilon$ = 1.0, k = 300, |ground| = 2,000), three proposal distributions, one held-out full-mixture pool. An \emph{oracle} ground set drawn from the full latent mixture --- the best a selector could do if conditioning could reach everything --- covers 0.387. A \emph{proposal-limited} ground set drawn from the base specs (which reach 25 of 60 modes) covers 0.212: a reachability gap of 17.4 points that no selection algorithm can close, because the deficit is in what the generator can be made to emit, not in how candidates are chosen. A \emph{refined-proposal} ground set (one spec per mode --- the target recursive refinement works toward) covers 0.374, recovering 92.5\% of the gap. This is Proposition 1 made quantitative: the (1 $-$ 1/e) guarantee binds relative to the proposal support, and conditioning --- not selection --- is where the missing coverage lives.

\textbf{6.7 Lattice size is not reachable dimension.} Using the sibling paper's slice-structured world (specs compose level directions additively; \texttt{../simulate\_slices.py}), we build two worlds with \emph{identical} lattice size $L^{A}$ = 1,024 distinct specs but different rank of the level-direction span: full-rank (rank 10, reachable dimension 13) versus low-rank (all axes forced into a shared 2-dim subspace: rank 2, reachable dimension 5). Same budget (n = 2,000 random-spec draws), same radius ($\varepsilon$ = 0.27): the low-rank world's own reachable pool is 68.1\% covered where the full-rank world's is 49.0\%. The number of distinct prompts you can write says nothing about how much there is to cover; covered volume is governed by reachable dimension, and a huge lattice over low-rank level directions is a thin space wearing a combinatorial costume. For coverage this distinction is sharper than for packing --- the denominator itself is set by reachable dimension --- and it is why our simulation worlds put a spec layer between policy and world rather than drawing from a global mixture: a simulator without conditioning structure cannot exhibit the effect and understates refinement (the sibling paper measured exactly this failure in its own first simulator).

\textbf{6.8 Refinement direction: density vs reach.} The sibling paper's negative result --- naive refinement hurt, because tight children near saturated parents concentrate mass where the corpus already sits --- ports to coverage, with a measurement subtlety worth recording. Nearest-neighbor distance CANNOT distinguish density-refinement from reach-refinement: in our first version of this check, child draws sat $\approx$ 1.04 from the nearest corpus point whether they had moved in-span or off-span, because a 1,500-item corpus is sparse in its own 13-dim occupied span. Density-vs-reach is a \emph{span} property. Measured with span-level metrics: \emph{local} children (tight, 0.25$\times$ displacement --- the naive move) produce draws of which 36.7\% are already covered by the pre-refinement corpus at the coverage radius, and add +0.02 effective dimensions --- almost pure density. Full-length \emph{isotropic} displacement produces 0\% pre-covered draws and +0.13 effective dimensions; \emph{transversality-guided} displacement 0\% and +0.14, with off-span displacement energy 0.915 vs 0.823 unguided. The honest reading: (i) the damning comparison is local-vs-far, and ``split the saturated cell into tighter sub-levels'' is only worth its cost if the children actually move; (ii) guidance beats unguided displacement, but by little \emph{here}, because this corpus occupies only a thin span (6 of 64 dimensions at the 95\% energy level) and a random direction is already 82\% transverse to it --- guidance matters most when the corpus has already spread, which is exactly the late-run regime where refinement fires.

\section{The honest denominator}
Recursive refinement creates an accounting problem that a coverage paper must not paper over: \emph{refinement changes the measure being covered}. The refine policy's single run admits three defensible coverage numbers (mid radius, redteam / persona):

\begin{table}[htbp]
\centering
\small
\begin{tabular}{lr}
\toprule
\textbf{denominator} & \textbf{covered fraction} \\
\midrule
base reachable pool (what the generator could do before refinement) & 0.492 / 0.501 \\
the policy's own final reachable pool (base + everything refinement opened) & 0.960 / 0.971 \\
the full latent space (oracle: all modes, reachable or not) & 0.381 / 0.538 \\
\bottomrule
\end{tabular}
\end{table}

Each number is true and each, alone, misleads. Against the \emph{fixed base pool}, refinement looks like a small regression (0.49 vs 0.54 for gated coverage without refinement): the budget it spent opening new territory was budget not spent filling old territory. Against the \emph{full latent space} it looks like the only policy that works (0.381 vs 0.214; 0.538 vs 0.272) --- but no base policy could have scored here at any budget, since 35 of 60 (20 of 40) modes were unreachable before refinement minted specs for them. And against its \emph{own final reachable pool} it looks nearly finished (0.96+), which is partly real (tight child specs are easy to cover) and partly the denominator being shaped by the same process being graded --- the refined proposal distribution concentrates exactly where the policy then samples.

\begin{figure}[htbp]
\centering
\includegraphics[width=\linewidth]{figures/fig_denominator.png}
\caption{Figure 5. One run, three denominators: the refine policy against the fixed base pool, its own final reachable pool, and the full latent space.}
\end{figure}

The operational face of the problem is the pool refresh of \S{}4: when refinement opens a region, the estimation pool must be re-drawn to include it, or the selector sees zero marginal gain precisely where the new territory is. We hit this live before we hit it in theory: in the pilot's first run, axis-conditioned candidates embedded $\approx$ 0.88 from an unconditioned-only pool with $\varepsilon _sel$ = 0.73 --- every candidate scored zero and the selector went blind --- and the fix was to build the reference pool from \emph{both} strata of the proposal process (\S{}8). The denominator must track the reachable distribution of the pipeline as it currently exists.

Our reporting rule, followed in every table above: state the denominator next to every coverage number, and report at least (i) the fixed pre-refinement pool, which is comparable across policies and cannot be gamed by growing the space, and (ii) the policy's own final reachable pool, labeled as self-referential, with the growth of the reachable set (modes reached, spec count) reported alongside. A single ``we covered X\%'' from a self-expanding pipeline is not a result; it is a choice of denominator.

\section{Live pilot}
\subsection{Setup and what it cost}
The pilot builds a red-team \emph{evaluation} prompt suite on the real stack: gpt-5.6-luna via OpenRouter for generation, judging, attractor mining and axis elicitation; local \texttt{nomic-embed-text} (768-d, unit-normalized) for embeddings. Severity is deliberately mild --- the items are benign user messages that probe whether an assistant adds disclaimers, resists sycophancy and leading questions, and stays calibrated --- and the judge scores a \texttt{benign} dimension that hard-gates acceptance. All state is append-only JSONL under \texttt{live\_logs/} and the run is resumable; it was in fact resumed mid-flight after a directory move, which is what the checkpointing is for.

The run reached its full target: \textbf{50/50 accepted items} in 50 steps (K = 3 candidates each), against a \textbf{50-item naive baseline} and a \textbf{360-item reference pool}. The pool has two strata --- 240 unconditioned one-shot draws and 120 axis-conditioned but \emph{unselected} draws --- for a reason discovered the hard way (\S{}8.3). Six axes were elicited, two refinements fired (8 axes final), and two ledger rounds mined attractors that became negative constraints. Judge quality averaged 8.89/10. Cost of the final instrumented segment: \textbf{74 calls, 36,056 prompt + 74,681 completion tokens, \$0.0968}, 394 s wall clock. Earlier segments (pool construction, the naive baseline, and the first 32 accepts) predate the usage ledger and are not included in that figure; total pilot spend was well under the \$3 budget but only the final segment is \emph{measured}, so we report that number rather than an estimate.

\subsection{Coverage results, and an honest negative}
Covered fraction of the 360-item pool at the three calibrated radii ($\delta$ = 0.592, $\varepsilon _sel$ = 0.770):

\begin{table}[htbp]
\centering
\small
\begin{tabular}{lrrr}
\toprule
\textbf{} & \textbf{$\varepsilon$ = 0.616} & \textbf{$\varepsilon$ = 0.770} & \textbf{$\varepsilon$ = 0.962} \\
\midrule
naive (n = 50) & 0.250 & \textbf{0.422} & \textbf{0.950} \\
method (n = 50) & 0.053 & 0.231 & 0.883 \\
\bottomrule
\end{tabular}
\end{table}

\textbf{On the pooled denominator the method loses, and we report that plainly.} The stratum breakdown at $\varepsilon _sel$ explains it and is the actual result:

\begin{table}[htbp]
\centering
\small
\begin{tabular}{lrr}
\toprule
\textbf{} & \textbf{unconditioned stratum (240)} & \textbf{conditioned stratum (120)} \\
\midrule
naive & \textbf{0.613} & 0.042 \\
method & 0.129 & \textbf{0.433} \\
\bottomrule
\end{tabular}
\end{table}

Each policy covers its own proposal distribution and barely touches the other's. The naive baseline is drawn from \emph{exactly} the prompt that produced the unconditioned stratum, so it enjoys a home-field advantage no selection rule can offset: it is covering the distribution it was sampled from. The axis-conditioned method lands somewhere else --- and covers \emph{that} region 10$\times$ better than naive does (0.433 vs 0.042). The two runs are not competing on one space; they are occupying different regions of it, and any single pooled number silently weights the comparison by how the pool was built (here 2:1 in the baseline's favor). This is the honest-denominator problem of \S{}7 appearing in live data, and it is sharper than the simulation version because here we cannot see the true reachable measure at all.

What the pilot therefore does and does not show. It does show: the machinery runs end-to-end on a real model at real cost; elicited axes move the generator into territory unconditioned sampling does not reach; coverage gains are measurable per item and decay as predicted. It does \textbf{not} show that the method beats naive sampling on coverage at n = 50 --- on the pooled denominator it does not, and we lack a neutral reference measure that would make the comparison fair. Building one requires sampling the generator's reachable space by a process independent of both policies, which at pilot scale we did not do. We flag this as the pilot's main methodological gap.

A second honest observation: \textbf{marginal gains hit zero and stayed there.} Coverage at $\varepsilon _sel$ plateaued at 0.208 from item \textasciitilde{}35 to item \textasciitilde{}47, with a small recovery to 0.231 after the second axis refinement fired at step 48. With 50 items and a 360-point pool, the method saturated what it could reach at that radius; the refinement-driven recovery at the very end is the mechanism of \S{}4 working, but 50 items is too few to see more than one such cycle.

\subsection{The blind-selector failure (why the pool has two strata)}
The first version of this pilot used a reference pool of unconditioned draws only. Every axis-conditioned candidate then scored a marginal gain of \emph{exactly zero}: measured post-hoc, method items sat at median distance 0.88 from the unconditioned pool while $\varepsilon _sel$ was 0.73, so no candidate's ball contained any pool point and the selector was choosing uniformly at random among ties for 30+ consecutive steps. The bug is conceptual, not numerical: if the reference measure does not include the region the proposal distribution actually reaches, coverage-greedy selection is blind exactly where the method is working. The fix --- pool the unconditioned and axis-conditioned strata --- is the live counterpart of the simulation's pool-refresh-after-refinement (\S{}4), and it is the single most transferable lesson of the pilot: \emph{the reference pool must be drawn from the same proposal process the policy uses, including every conditioning mechanism.}

\subsection{Literal vs latent diversity, and a kernel caveat}
Tracking both n-gram and embedding diversity as the corpus grows reproduces a divergence reported at much larger scale on this model family: the two move in \textbf{opposite directions}.

\begin{table}[htbp]
\centering
\small
\begin{tabular}{lrrrr}
\toprule
\textbf{n} & \textbf{distinct-1} & \textbf{distinct-2} & \textbf{Vendi (centered)} & \textbf{Vendi (uncentered)} \\
\midrule
5 & 0.614 & 0.944 & 3.98 & 3.62 \\
10 & 0.482 & 0.861 & 8.01 & 5.04 \\
20 & 0.367 & 0.809 & 16.16 & 7.37 \\
30 & 0.338 & 0.791 & 23.65 & 10.00 \\
40 & 0.312 & 0.772 & 30.55 & 11.87 \\
50 & 0.295 & 0.770 & 37.06 & 13.54 \\
\bottomrule
\end{tabular}
\end{table}

Literal diversity falls monotonically (distinct-2 0.944 $\to$ 0.770) while latent diversity rises monotonically (centered Vendi 3.98 $\to$ 37.06). Reporting either alone supports an opposite conclusion about the same corpus. The mechanism is plain enough --- a growing corpus reuses vocabulary and sentence frames while still spreading in meaning-space --- but it means ``diversity'' must be qualified by level, and a coverage objective defined on embeddings is silent about literal repetition. (Against the naive baseline the method is slightly \emph{less} literally diverse at distinct-1, 0.295 vs 0.341, and slightly more diverse latently, centered Vendi 37.06 vs 30.34, with markedly lower self-repetition, max 3-gram Jaccard 0.036 vs 0.083.)

\textbf{The uncentered Vendi is a kernel artifact.} Same-domain embeddings share a mean direction, and an uncentered linear kernel measures that shared cone as if it were content. On this data the two readings differ by 2.7$\times$ (13.54 vs 37.06) and, worse, they order the growth curve differently in slope. We report both everywhere and treat only \emph{relative} comparisons at fixed n as meaningful for the uncentered form.

\textbf{Is $\varepsilon$ measuring the cone or the content?} This is the question the caveat raises for every coverage number computed with cosine distances, so we measured it (\texttt{cone\_diagnostics} in \texttt{summary\_live.json}). Our pool's mean pairwise cosine is \textbf{0.444} (p05 0.330, p95 0.591) --- a wide cone, not the \textasciitilde{}0.9 concentration seen in some same-domain embedding sets --- with mean direction norm 0.667. Our $\varepsilon _sel$ = 0.770 corresponds to cosine similarity 0.704, which sits at the \textbf{1.8th percentile} of pairwise distances: an $\varepsilon$-ball at this radius captures genuine near-duplicates, not the bulk of the cone. So for this pilot the coverage numbers are measuring content. We record the diagnostic because the answer is embedder- and domain-dependent: with a tighter cone the same calibration rule would have produced an $\varepsilon$ inside the cone's width, and the resulting ``coverage'' would have been a restatement of the embedding geometry. Any coverage result computed on cosine distances should publish this check.

\begin{figure}[htbp]
\centering
\includegraphics[width=\linewidth]{figures/fig_live_pilot.png}
\caption{Figure 6. Live pilot on gpt-5.6-luna (n = 50). Left: coverage growth of the method run versus the naive baseline at three radii on the pooled denominator, where naive leads. Center: per-item marginal gains, showing the plateau from item \textasciitilde{}35 and the recovery after the refinement at step 48. Right: the same runs split by pool stratum -- each policy covers its own proposal distribution, which is what the pooled number obscures.}
\end{figure}

\begin{figure}[htbp]
\centering
\includegraphics[width=\linewidth]{figures/fig_literal_vs_latent.png}
\caption{Figure 7. Left: literal (n-gram) and latent (embedding) diversity move in opposite directions as the corpus grows. Right: uncentered versus mean-centered Vendi on the same sets -- the shared mean direction compresses the uncentered score by roughly 2.7x, so only relative comparisons at fixed n are trustworthy.}
\end{figure}

\section{Diagnosing a negative result: three ways to measure coverage wrong}
Section 8 reported that our live pilot lost to a naive baseline. Rather than accept that, we treated it as a bug report and worked through it. The diagnosis took three measurement fixes and one conceptual one, and the corrected comparison reverses the result. We report the whole sequence, because each wrong version is a mistake a coverage paper can easily make.

\subsection{What we first measured}
The parent project's \texttt{real/} directory contains two 10,000-item corpora of real gpt-5.6-luna output produced by plain repeated prompting (\texttt{dalle\_naive}, \texttt{psychometric\_naive}), with saved embeddings. No selection policy shaped them, so they look like an ideal policy-independent reference measure, and we used one as the denominator for every arm of the parent project's competitive matrix (\texttt{real\_coverage.py}). The result was stark: at the calibrated mid radius, coverage was \textbf{anti-correlated with every diversity metric}. In the DALL-E domain, naive scored 0.137 and high-temperature 0.135, while Self-Instruct scored 0.015, Evol-Instruct 0.011, and Persona-Hub, the parent method, and its vision-steered variant all scored \textbf{0.000} --- the same ordering, reversed, as distinct-2 (0.284 for naive rising to 0.699 for the vision arm) and centered Vendi (50.3 rising to 92.0).

A number that ranks methods in exact reverse order of every other diversity measure is not measuring what its name says. Three things were wrong.

\subsection{Confound 1: the denominator encoded a distribution, not a space}
Unconditioned prompting is not ``the reachable space''; it is one region of it. Scoring conditioned generation against it asks ``how much of what the model would say anyway does this method reproduce?'' --- and a method whose entire purpose is to leave that region must score low. This is the honest-denominator problem of \S{}7 in its sharpest form: the denominator was not neutral, it was the naive policy's own distribution.

Our first repair --- a union of all arms --- was not enough, because an unweighted union is dominated by whichever stratum has the most points (the 5,000-point naive reference), smuggling the same distribution back in. The fix is a \textbf{stratum-balanced} reference in which every generation method contributes exactly the same number of points, scored \textbf{leave-one-out} so no corpus is ever scored against itself.

\subsection{Confound 2: epsilon was calibrated on the wrong set}
We calibrated $\varepsilon$ on the dense naive corpus and then applied it to other denominators. But $\varepsilon$ must be commensurate with the reference measure it is applied to. A radius tuned to the tightest region of the space makes every diffuse method score zero \emph{everywhere}: the conditioned arms have median within-corpus nearest-neighbor distances of 0.50--0.64, while $\varepsilon$ was 0.253. The balls were smaller than the typical spacing of the sets being measured, so they caught nothing. Calibrating $\varepsilon$ on the reference actually in use raises it to 0.467 (DALL-E) and 0.458 (psychometric), and we now report where $\varepsilon$ sits in that reference's own pairwise-distance distribution --- the 4.1st and 0.9th percentiles respectively, confirming the balls capture near-neighbors rather than the bulk of the embedding cone.

\subsection{Confound 3: comparing pipelines conflates generation with selection}
``Persona-Hub vs ours'' bundles two separable questions: what candidates get generated, and which get kept. This paper's contribution is a selection rule. Selection is only comparable when every rule sees the same candidates.

We also verified the estimator itself before trusting it further (\texttt{verify\_estimator.py}): against the analytic volume of a ball in a uniform cube it is within 3 Monte-Carlo standard errors at d = 2, 3, 4; the union of nine disjoint balls matches m$\cdot$V to within 2.0e-4; duplicate centers add exactly zero coverage (the submodularity-critical property a naive sum-of-counts implementation would violate); and on a 1-D instance with an analytically known optimum, greedy attains it exactly (70/70 points). The estimator was not the problem.

\subsection{The corrected comparison: budget-matched selection}
\texttt{real\_benchmark.py} fixes all three. Every method selects k = 400 items from one shared candidate pool, so all consume exactly the same number of generations and differ only in which items they keep --- the comparison is budget-matched by construction, not merely n-matched. The reference is stratum-balanced; $\varepsilon$ is calibrated on it.

\textbf{A fourth error, found by auditing our own win.} Our first version of this benchmark built one incidence matrix against the reference, handed it to the selectors, and then scored against that same reference. Coverage-greedy was therefore optimizing the evaluation set directly, and its margin measured tautology rather than generalization. We split the reference into two disjoint halves: an estimation half the selector may see (standing for a pipeline's own cheap draws) and a held-out half used only for scoring. The win survives, and it shrinks: from +0.039/+0.107 under leakage to \textbf{+0.020/+0.033} below. We report the corrected numbers and flag the leaky ones as what they were.

Five real baselines from the literature, implemented faithfully: farthest-point traversal (Gonzalez 1985), SemDeDup (Abbas et al. 2023), Self-Instruct's ROUGE-L $\leq$ 0.7 sequential filter (Wang et al. 2023), greedy MAP-DPP (Kulesza \& Taskar 2012; Chen et al. 2018), and a uniform random draw (what a corpus yields as shipped).

\textbf{DALL-E instructions} (pool 700, reference 700, k = 400, $\varepsilon$ = 0.467):

\begin{table}[htbp]
\centering
\small
\begin{tabular}{lrrrr}
\toprule
\textbf{selector} & \textbf{covered @$\varepsilon$} & \textbf{min-gap} & \textbf{Vendi (centered)} & \textbf{distinct-2} \\
\midrule
\textbf{coverage-greedy (ours)} & \textbf{0.4829} & 0.299 & 106.5 & 0.677 \\
random & 0.4629 & 0.198 & 87.2 & 0.580 \\
ROUGE-L filter (Self-Instruct) & 0.4629 & 0.208 & 87.7 & 0.583 \\
SemDeDup & 0.4600 & 0.336 & 96.2 & 0.618 \\
coverage-stream, K=8 (ours) & 0.4200 & 0.283 & 47.7 & 0.723 \\
k-center (Gonzalez) & 0.4114 & \textbf{0.491} & 111.7 & 0.679 \\
greedy MAP-DPP & 0.3829 & 0.402 & 113.6 & 0.683 \\
\bottomrule
\end{tabular}
\end{table}

\textbf{Psychometric items} (pool 1260, reference 1266, k = 400, $\varepsilon$ = 0.458):

\begin{table}[htbp]
\centering
\small
\begin{tabular}{lrrrr}
\toprule
\textbf{selector} & \textbf{covered @$\varepsilon$} & \textbf{min-gap} & \textbf{Vendi (centered)} & \textbf{distinct-2} \\
\midrule
\textbf{coverage-greedy (ours)} & \textbf{0.3901} & 0.036 & 87.8 & 0.519 \\
coverage-stream, K=8 (ours) & 0.3785 & 0.072 & 56.8 & 0.496 \\
SemDeDup & 0.3575 & 0.234 & 94.5 & 0.532 \\
ROUGE-L filter (Self-Instruct) & 0.3575 & 0.080 & 95.6 & 0.537 \\
random & 0.3533 & 0.000 & 77.3 & 0.510 \\
greedy MAP-DPP & 0.0988 & 0.533 & 136.0 & 0.575 \\
k-center (Gonzalez) & 0.0557 & \textbf{0.748} & 135.2 & 0.560 \\
\bottomrule
\end{tabular}
\end{table}

Coverage-greedy wins both domains on held-out reference points, beating the best baseline by +0.020 (1.043x) on DALL-E and +0.033 (1.091x) on psychometric items. These are modest margins, and they are the honest ones: the leaky version of this same experiment reported 1.08x and 1.29x. The deployable streaming variant, which never rescans the pool, is mid-field rather than second --- with the estimation reference halved it has fewer points per ball to work with, and the best-of-K rule is more sensitive to that than full greedy is.

\begin{figure}[htbp]
\centering
\includegraphics[width=\linewidth]{figures/fig_benchmark.png}
\caption{Figure 8. Budget-matched selection on real gpt-5.6-luna corpora: every method picks k = 400 items from one shared candidate pool, scored against a stratum-balanced leave-one-out reference. Left and center: covered fraction by selector, ours in blue. Right: the covering/packing dissociation across both domains -- k-center maximizes min-gap and minimizes coverage.}
\end{figure}

Two things in these tables matter more than the headline. First, the \textbf{covering/packing dissociation is far more extreme on real data than in simulation}: k-center wins min-gap by a wide margin in both domains (0.491 and 0.748) and finishes \emph{last} on coverage in the psychometric domain at 0.056 --- a sixth of what a uniform random draw achieves. Choosing the k-center objective for a covering problem is not a small mistake. Second, \textbf{high spectral diversity does not imply coverage}: MAP-DPP and k-center post the highest Vendi scores (114--136) and the lowest covered fractions, because determinant and min-gap both reward pushing to the boundary, while our streaming variant posts the lowest Vendi in the DALL-E table (47.7) at a covered fraction near the middle. Vendi and coverage are different quantities, and a diversity claim has to name which one it means.

\subsection{An implementation trap worth naming}
Twice during this work a run silently degraded rather than failing, in the same way, and the failure is easy to mistake for a modelling result. This model bills \emph{reasoning} tokens against the same \texttt{max\_tokens} budget as the visible answer. A limit sized for the answer alone therefore returns an \textbf{empty content field} once the prompt gets harder --- and a prompt gets harder exactly when our method is working, because the attractor ledger keeps appending ``do not do X'' constraints. The symptom is distinctive: cost keeps accruing at the normal per-call rate while accepted-item yield collapses (in our instruction run, from 48 accepted items per 48-call batch to 2, at unchanged cost per batch). Read naively, that looks like the ledger suppressing the generator's ability to produce valid output --- a finding about negative constraints. It is an off-by-a-token-budget bug. Raising the limit from 400 to 1,500 restored full yield immediately. Any pipeline that adds accumulating constraints to its prompts should size token budgets for reasoning plus answer, and should alert on yield-per-call, not just on errors.

\subsection{Duplication sets the scale, and conditioning removes it}
The psychometric corpora make the case for deduplicating before measuring. Naive prompting produced a \textbf{73.6\% exact-duplicate rate} over 10,000 generations --- 2,645 unique texts, with a single item appearing \textbf{2,726 times}. High-temperature sampling barely helped (71.0\%). Conditioned arms essentially eliminated it: Self-Instruct 4.7\%, Evol-Instruct 1.2\%, Persona-Hub 0.7\%, the parent method 0.0\%. In the DALL-E domain every arm sits at 0.0\%, so this is a property of the domain's answer space, not of the model in general.

Two consequences. First, any coverage number computed over a corpus like that is measuring duplication: against the raw naive pool the naive arm scores 0.840, and against the deduplicated pool the same items score 0.350. We report deduplicated coverage as the primary number throughout. Second, $\varepsilon$-calibration by nearest-neighbor distance breaks down completely --- the median nearest-neighbor distance in the raw psychometric pool is \textbf{exactly 0.0000}, because the median item has an identical twin. Duplication does not just bias distance-based diversity metrics, it can make them undefined. And since temperature does not fix it while conditioning does, this is direct evidence for the paper's thesis that conditioning, not sampling temperature, is what buys coverage.

\subsection{Head-to-head against released corpora, and why the coverage metric fails across corpora}
The comparisons so far pit our selector against reimplementations. The sharper question is whether our corpus holds up against the artifacts people actually use, so we compare against three released corpora as downloaded: \textbf{Alpaca} (Self-Instruct, 52,002 items), \textbf{PersonaHub} (50,000), and \textbf{WizardLM Evol-Instruct V2} (143,000). Against them we place three arms generated from gpt-5.6-luna at 450 items each: unconditioned prompting, the same prompt at T = 1.6, and axis-conditioned generation with an attractor ledger (\texttt{instruction\_gen.py}, \$0.376, 564 s). Matched n = 60 after deduplication, identical embedder, instruction field only.

\textbf{The coverage metric does not survive this comparison, and that is the result.} Ranked by covered fraction of a stratum-balanced reference at the calibrated radius, with each corpus's own spread alongside:

\begin{table}[htbp]
\centering
\small
\resizebox{\linewidth}{!}{%
\begin{tabular}{lrrrrr}
\toprule
\textbf{source} & \textbf{covered @$\varepsilon$} & \textbf{median self-NN} & \textbf{Vendi (centered)} & \textbf{distinct-2} & \textbf{dup rate} \\
\midrule
ours: naive & 0.3125 & 0.280 & 16.3 & 0.393 & 0.198 \\
ours: high-temp T=1.6 & 0.3063 & 0.305 & 16.8 & 0.402 & 0.207 \\
ours: axis-conditioned & 0.1271 & 0.696 & 40.3 & 0.738 & 0.000 \\
WizardLM Evol-Instruct (143k) & 0.0521 & 0.915 & 47.6 & 0.759 & 0.010 \\
Alpaca (Self-Instruct, 52k) & 0.0083 & 0.902 & 47.3 & 0.868 & 0.000 \\
PersonaHub (50k) & 0.0063 & 0.901 & 47.3 & 0.821 & 0.000 \\
\bottomrule
\end{tabular}
}
\end{table}

The ordering is exactly inverted. Across these six corpora, covered fraction correlates \textbf{$-$0.991} with median within-corpus nearest-neighbor distance and \textbf{$-$0.986} with centered Vendi. The metric ranks the \emph{most degenerate} corpus first: our naive arm, which repeats itself 19.8\% of the time and has a within-corpus spacing of 0.280, ``covers'' 50x more than PersonaHub.

The mechanism is now clear, and it subsumes every confusion earlier in this section. Coverage at a fixed $\varepsilon$ rewards corpora whose own items are packed tightly \emph{relative to $\varepsilon$}. A tight corpus self-covers: its 80 contributed reference points all sit inside a few balls, so a handful of selected items sweeps its entire stratum, and any other stratum occupying the same region comes free (our naive and high-temp arms overlap, which is why both sit near 0.31 $\approx$ 2/6). A spread corpus cannot even cover its own stratum, because its typical spacing (0.90) exceeds $\varepsilon$ (0.75). \textbf{No single $\varepsilon$ can be simultaneously appropriate for a corpus with spacing 0.28 and one with spacing 0.90.} Cross-corpus coverage comparison at a fixed radius is not a valid measurement, and tuning $\varepsilon$ per corpus until the preferred method wins would be worse than not measuring at all.

This is the honest-denominator thesis (\S{}7) in its strongest form. Coverage is a ratio, and here both the numerator's scale and the denominator's composition are set by the corpora being compared. The metric is well-posed \emph{within} a shared pool --- which is exactly the condition \S{}9.5 enforces, and why the selection benchmark there is the defensible experiment --- and ill-posed across corpora with different intrinsic scales.

\textbf{What the scale-free metrics say.} Diversity measures that need no shared radius do support a clear conclusion, and it is a mixed one for us. Axis conditioning is transformative \emph{relative to our own baselines}: it raises centered Vendi from 16.3 to 40.3 (2.5x), distinct-2 from 0.393 to 0.738 (1.9x), within-corpus spacing from 0.280 to 0.696, and drives exact duplication from 19.8\% to 0.0\%. High temperature does essentially nothing on any of these (Vendi 16.8, distinct-2 0.402, duplication 20.7\%) --- a third independent replication, on a third domain and our own generator, of the finding that conditioning rather than temperature is what buys diversity.

But conditioning does \textbf{not} close the gap to the released corpora, which sit at Vendi 47.3--47.6 and spacing \textasciitilde{}0.90 against our 40.3 and 0.696. Our corpus is 450 items from one small model; theirs are 50k--143k items from different, mostly stronger generators, with human curation in at least one pipeline. We reach roughly 85\% of their centered-Vendi at 1/100th the scale, which we think is a real result for the method and is not the same as beating them. No framing fixes the part we did not win.

\section{Related work}
\textbf{Submodular maximization.} The (1 $-$ 1/e) greedy guarantee for monotone submodular objectives is Nemhauser, Wolsey \& Fisher (1978); hardness of beating it for max-coverage is Feige (1998). Lazy greedy (Minoux 1978) and stochastic greedy (Mirzasoleiman et al. 2015) make greedy cheap; sieve-streaming (Badanidiyuru et al. 2014) gives the single-pass regime our sequential setting approximates. Facility location --- of which our pool objective is the thresholded special case --- is the standard submodular model of ``representativeness'' in data summarization (Lin \& Bilmes 2011).

\textbf{Packing vs covering, k-center vs k-median.} Farthest-point traversal (Gonzalez 1985) 2-approximates k-center --- the packing objective the sibling paper streams --- while k-median/facility-location minimizes \emph{average} service distance and behaves like our coverage objective (measure-weighted, bulk-first, outlier-indifferent). Coreset constructions (Har-Peled \& Mazumdar 2004) make the same distinction: an $\varepsilon$-coreset for k-median weights by mass, a k-center net does not. Our contribution is not new theory for these objects but the identification of \emph{which one the corpus-construction product requirement actually is}, plus measurement of what optimizing the wrong one costs.

\textbf{Quality-diversity search.} MAP-Elites (Mouret \& Clune 2015) maintains an ``illuminated map'': one elite per cell of a hand-designed behavior space --- equal allocation over a fixed grid, in our vocabulary, with quality as the within-cell objective. Our method differs in that cells are elicited from the generator and refined recursively rather than fixed, allocation is marginal-gain-driven rather than one-per-cell, and coverage is measured against a reachable distribution rather than a designer-set grid; \S{}6.5's equal-allocation row quantifies what the fixed-grid choice costs at small $\varepsilon$.

\textbf{Diversity metrics and dedup.} The Vendi score (Friedman \& Dieng 2022) is our spectrum-level diversity readout, computed at O($nD^{2}$ + $D^{3}$) via the Gram trick as in the sibling paper. SemDeDup (Abbas et al. 2023) removes semantic near-duplicates from web-scale corpora --- the deletion-side mirror of our selection problem: both spend a budget to maximize non-redundant semantic territory. DPPs (Kulesza \& Taskar 2012) give probabilistic diverse selection; exact DPP sampling scales poorly in n and optimizes a determinant (volume) rather than covered measure, though greedy MAP-DPP behaves similarly to facility-location greedy in practice.

\textbf{Eval-suite construction.} Perez et al. (2022) generate evaluations with LMs at scale; red-teaming work (Ganguli et al. 2022) documents exactly the clustering-into-attack-families failure that motivates a coverage objective for safety suites.

\textbf{Synthetic instruction corpora, and the baselines we compare against.} Self-Instruct (Wang et al. 2023) bootstraps instructions from a seed pool with few-shot exemplars and a ROUGE-L similarity filter; the released Stanford Alpaca corpus (Taori et al. 2023, 52,002 items) is its canonical output. Evol-Instruct / WizardLM (Xu et al. 2023) instead \emph{mutates} existing instructions with depth and breadth evolution operators, released as WizardLM\_evol\_instruct\_V2 (143,000 items). Persona-Hub (Ge et al. 2024) conditions generation on a catalogue of \textasciitilde{}1B personas, releasing 50,000 persona-synthesized instructions; it is the closest published relative of latent-axis conditioning, differing in that its catalogue is mined at scale and flat, where our axes are elicited from the generator and refined into a tree. AttrPrompt (Yu et al. 2023) similarly conditions on attribute dimensions. We compare against the released artifacts of the first three (\S{}9.7) rather than against reimplementations, since a reimplementation can be weak in ways that flatter us. On the selection side our baselines are the standard ones for diverse subset choice: SemDeDup (Abbas et al. 2023), farthest-point traversal (Gonzalez 1985), and greedy MAP inference for DPPs (Kulesza \& Taskar 2012; Chen, Zhang \& Zhou 2018).

\section{Relation to the infinite-horizon paper}
The two papers are a packing/covering dual pair, and the duality is the k-center/k-median one:

\begin{table}[htbp]
\centering
\small
\begin{tabular}{lrr}
\toprule
\textbf{} & \textbf{infinite horizon (sibling)} & \textbf{finite budget (this paper)} \\
\midrule
horizon & unbounded stream & known n = 10,000 \\
objective & max-min gap + typicality anchor & $\varepsilon$-ball covered measure of $\mu$ \\
classical problem & k-center / packing & max coverage / facility location \\
optimum's habitat & boundary, extremes & bulk, measure-weighted \\
junk incentive & attracted (needs gate badly) & indifferent (gate still needed for quality) \\
planning & impossible (no n) & allocation across cells (\S{}3.6) \\
guarantee & none exact (packing is 2-approx streaming) & (1 $-$ 1/e), $1/2$ streaming \\
refinement's role & capacity so the stream never exhausts & reach --- and a denominator liability \\
\bottomrule
\end{tabular}
\end{table}

The method stack transfers almost unchanged --- elicited axes, judge gate, ledger, recursive refinement are all objective-agnostic scaffolding for steering a generator; only the selection rule and its state differ (running second-moment matrix and min-gap scales there; reference pool and covered mask here). The sibling's slice-world results are directly load-bearing for this side of the duality. At an identical budget of 10,000 generations, its measured policy ladder runs from fixed-prompt (Vendi 1.06, reachable dimension 13) through random-spec (6.70) and calculus-guided specs (7.93) to calculus plus transversality-guided refinement (Vendi 11.23, reachable dimension 24) --- a 10.6$\times$ diversity gain from conditioning and refinement alone, with refinement raising the reachable dimension from 13 to 24 (transversality theorem P4 in action). Two of its negative results transfer with particular force here. First, \emph{naive refinement hurt}: children split isotropically near their saturated parents concentrated probability mass exactly where the corpus already sat; refinement pays only when the new directions are transverse to the occupied span --- our \S{}6.8 reproduces this under the coverage objective (density, not reach). Second, its \emph{inability-to-be-novel} metric was best ($-$0.104) for the fixed-prompt policy precisely because that policy covered almost nothing and so consumed no headroom. The coverage-side analogue of that trap: a method that concentrates every $\varepsilon$-ball in one small region leaves the rest of the space ``available'', which flatters any residual- or headroom-based metric. This is why every table in \S{}5 reports covered-fraction \emph{jointly} with a within-corpus diversity measure (Vendi) and why coverage is always measured against the reachable manifold's measure, never against what happens to be left over. Its junk-conflation measurement is also starker than ours and worth quoting for eval-suite builders: in a hard-floor novelty bank with no typicality gate, the junk fraction of accepted items rose from 7.9\% in the first half of the run to 48.2\% in the second --- once legitimate space saturates, the only candidates still clearing a novelty constraint are the incoherent ones. A red-team suite or persona corpus built that way looks maximally diverse by every embedding metric and is half garbage, silently.

The deepest difference is epistemic: the infinite-horizon paper never has to say what fraction of anything it achieved, while a coverage claim is \emph{always} a ratio, which is why the denominator problem (\S{}7) exists only on this side of the duality. Conversely, the budgeted side gets something the unbounded side cannot have: a completion semantics. A coverage run can report ``we are 62\% done at resolution $\varepsilon$ and the marginal item now buys 0.1\% '' --- a spend/stop signal no packing objective provides.

\section{Beating Alpaca: a diagnosis-driven campaign}
Earlier sections reported that our corpora lost to the released instruction sets on scale-free coverage of human-written Dolly instructions --- Alpaca at 0.365 against our best 0.103 --- and were told to treat that as a bug report rather than a verdict. This section is the debugging session, run to completion. It ends with a corpus that takes first place among twelve, beating Alpaca 0.4441 to 0.3722 on a reference no part of the pipeline ever aimed at, at matched evaluated-\emph{n}, with the highest precision in the field. Every step was measured before it was believed, three hypotheses died on the way, and the mechanism that finally worked is a direct instantiation of the covering theory of \S{}3 --- including one component that \emph{hurt} until the objective's own logic said to invert it.

\subsection{Making the fight fair}
Alpaca is synthetic, but it is not \emph{unseeded}: Self-Instruct bootstraps from 175 human-written seed tasks and few-shot-prompts every generation from them, so its distribution is anchored to human instruction style by construction. Scoring a seedless pipeline against human-distribution coverage measures the head start, not the method. The controlled design: every arm receives the same 175 seed tasks (the actual \texttt{seed\_tasks.jsonl} Alpaca grew from), the same generator, the same budget counted in \textbf{generator calls} --- so rejection and selection losses are paid, not hidden --- the same embedder, and evaluation on a held-out half of Dolly that nothing in any loop ever reads. The steering half (STEER) and the evaluation half are disjoint; the generator never sees one word of either.

\subsection{Three hypotheses, two deaths}
\textbf{Reservoir depth --- dead.} Alpaca evaluates as 450 draws from 52,002 items; our arms held \textasciitilde{}2,400. Constraining Alpaca to a random 2,400-item reservoir left its score unchanged (0.3743). Its edge is per-item, not statistical depth.

\textbf{Register offset --- dead, and inverted.} Our few-shot items were \emph{closer} to Dolly per-item (median nearest-reference distance 0.307 vs Alpaca's 0.355) at comparable length. Proximity was never the problem.

\textbf{Within-corpus twinning --- confirmed, massively.} In a 450-item sample of our few-shot corpus, \textbf{80.2\% of items have a near-twin} at cosine distance < 0.15 (median within-sample nearest-neighbour distance 0.041, with exact duplicates at the 10th percentile). Alpaca: 0.9\%. Dolly itself: 1.3\%. The generator that produced a 72.3\%-duplicate psychometric bank produces, under few-shot, a corpus of near-paraphrases that \emph{looks} varied. Our items each sit close to some reference point; they sit close to the \emph{same} reference points.

\subsection{The mechanism, assembled from the theory}
Four components, each an answer to a measured failure:

\begin{enumerate}
  \item \textbf{Aim by retrieval} (the inverse-oracle workaround). We cannot decode a target embedding into text, but we can sample an \emph{uncovered} reference region and retrieve the nearest texts we already own --- seeds plus our own corpus --- as the few-shot exemplars. Few-shot does the decoding; generation lands where nothing is. The reference enters as embeddings on the steering side only.
  \item \textbf{Radius-adaptive mode choice} (facility location, enacted in the prompt). Each reference point carries its own k-NN radius. A \emph{tight} radius means a dense region: coverage there demands typicality, so the prompt switches to imitate-closely mode --- no diversification contracts at all. A \emph{wide} radius means sparse territory: full diversification with explicit in-prompt negatives (the nearest own outputs to the target, shown as ``not like these''). Targets are sampled density-weighted ($\propto$ 1/$r^{2}$), because dense cores are many-points-per-item and expire fast. This single change moved the interim score from 0.3014 to 0.4485.
  \item \textbf{Twin suppression in the literal channel}: level-usage balancing (the coverage calculus's headroom term applied at conditioning time), a terse-register mandate, and the in-prompt negatives.
  \item \textbf{Keep everything.} Coverage is monotone; at a generation-matched budget every discard is a tax no selector can repay.
\end{enumerate}

\subsection{The ladder: every rung falsifies a hypothesis}
All arms, full budget, scored on the held-out half (full-corpus AUC over k $\in$ \{3,5,10,20\}):

\begin{table}[htbp]
\centering
\small
\begin{tabular}{lrrr}
\toprule
\textbf{arm} & \textbf{kept} & \textbf{AUC} & \textbf{what it isolates} \\
\midrule
\textbf{v4: aim + density-adapt + keep-all} & 2,398 & \textbf{0.6597} & the assembled mechanism \\
Self-Instruct (reimplemented) & 1,833 & 0.3262 & +ROUGE filter $\approx$ nothing (+0.0016) \\
bare few-shot from seeds & 2,399 & 0.3246 & seeds do the anchoring; corpus is a twin farm \\
axis conditioning, keep-all & 1,310 & 0.3147 & conditioning alone under-performs few-shot \\
\textbf{orthogonalized conditioning (v3)} & 1,152 & \textbf{0.2941} & \textbf{orthogonalization is anti-coverage} \\
selective 1-of-8 & 304 & 0.2591 & selection is a coverage tax \\
\bottomrule
\end{tabular}
\end{table}

Two rungs deserve emphasis. The celebrated ROUGE filter of Self-Instruct contributes 0.0016 --- the seeds contribute essentially everything. And the orthogonalized-conditioning arm --- our own Part-I machinery, applied faithfully --- scored \emph{below plain conditioning}. That is not an implementation failure; it is the theory's prediction landing on our own method. Orthogonalization pushes generation away from the occupied span, which is away from the reference-dense core that coverage is paid to fill: the k-center mistake, committed at the conditioning level. The packing tool, applied to the covering objective, loses to doing nothing.

\subsection{The verdict}
Scored on the STEER-disjoint reference half --- no corpus under test, ours or theirs, was aimed at any scored point --- at matched evaluated-\emph{n} = 450, identical estimator (Naeem et al. coverage/density with reference-side k-NN radii, AUC over k):

\begin{table}[htbp]
\centering
\small
\begin{tabular}{lrrr}
\toprule
\textbf{rank} & \textbf{corpus} & \textbf{AUC} & \textbf{precision} \\
\midrule
1 & \textbf{ours v4} & \textbf{0.4441} & \textbf{0.973} \\
2 & Alpaca (52k) & 0.3722 & 0.87 \\
3 & ours selective & 0.2591 & 0.93 \\
4 & PersonaHub (50k) & 0.2532 & 0.78 \\
5 & WizardLM Evol-Instruct (143k) & 0.2329 & 0.77 \\
6--9 & fair-fight ablation arms & 0.2102--0.2188 &  \\
10--12 & unseeded arms (naive, high-temp, axis-only) & 0.095--0.101 &  \\
\bottomrule
\end{tabular}
\end{table}

A 19\% margin over Alpaca, achieved at 1/20th of Alpaca's generation budget, with the field's highest precision. The interim leak-free check told the same story mid-run (0.4276 vs 0.3666 at n=528 of budget), so the result does not depend on the final corpus size.

\textbf{The asymmetry, stated plainly.} v4 consumes a sample of the target distribution --- as embeddings, on the steering side only. Alpaca had no such input. That is not an incidental advantage; it is the method's entire point --- coverage is always coverage \emph{of} something, and a method that may specify the something should. The like-for-like no-reference comparison is the ladder's middle rungs, where our no-STEER arms sit at parity with (slightly below) the Self-Instruct family. The claim this section supports is therefore precise: \emph{given a specification of the space to cover --- even one the generator itself never reads --- retrieval-aimed, density-adaptive conditioning covers it far better than any tested method covers it blind, and better than the strongest seeded baseline covers it at twenty times the budget.}

\textbf{The dissociation's sharpest exhibit.} v4's mean-centered Vendi is 62.1 --- the \emph{lowest} of any fair-fight arm (its density: 1.69, nearly twice uniform). The corpus that wins coverage is the internal-diversity loser, deliberately: it spends items where the reference measure is, twins included. Anyone selecting a ``diverse synthetic corpus'' by Vendi score would rank the winner last.

\subsection{What it cost, and what broke}
The campaign consumed roughly \$6 of generation across six arms. Two negative results were kept: hypothesis deaths H1 and H2 above, and the v3 inversion. One retraction was made along the way: an early probe reported plan-compliance and alignment figures computed on corrupted prompt-audio pairings in a different domain (\S{}audio); the repaired figures replaced them and the incident is registered in the provenance ledger. Every number in this section carries a P3 provenance tag; the arms' logs, the seed file, both Dolly halves, and the evaluation harness are in the repository.

\section{Limitations}
\begin{itemize}
  \item \textbf{Simulation realism.} Gaussian mixture worlds with a spec layer capture mode structure, skewed mass, junk, and limited base reach, but not the ways real embedding geometry misrepresents semantic distinctness ( anisotropy, hubness), nor judges whose errors correlate with novelty --- the live pilot mitigates but at n = 50, not 10,000.
  \item \textbf{$\varepsilon$ is chosen, not learned.} All guarantees are per-$\varepsilon$; our calibration (quantiles of reachable NN distance) is a heuristic, and \S{}6.5 shows the cost of optimizing at the wrong radius. A multi-resolution objective (integrating coverage over an $\varepsilon$ prior) is the obvious next step.
  \item \textbf{The pool is the measure.} Everything is relative to the reachable pool; behaviors the base generator cannot emit are invisible until refinement opens them (\S{}7), and pool refresh is only as good as the refinement trigger.
  \item \textbf{Selection bias on the estimation pool.} We hold out evaluation pools in simulation, but the live pilot's selection and reporting share one 240-item pool (with the bias direction stated); a larger live run should split them.
  \item \textbf{Live scale.} n = 50, one seed, one model; the live pilot is evidence of mechanism, not of effect size at n = 10,000.
  \item \textbf{What the benchmark win is and is not.} Section 9.5 shows our \emph{selection rule} beating five baselines on a shared candidate pool. That is a claim about selection, not about generation: the pool was built by other methods, and a selector cannot cover territory its pool never reached (\S{}6.6). The reachability gap remains the binding constraint, and closing it is a generation problem.
  \item \textbf{Single embedder, single similarity geometry.} Every coverage number here is computed in \texttt{nomic-embed-text} space with cosine distance. The parent project measured a correlation of only r = 0.155 between text-embedding and CLIP-image-embedding pairwise similarity on rendered DALL-E outputs, which means text-space coverage is close to blind to whether the \emph{rendered artifacts} are diverse. For any domain with a downstream rendering step, the objective should be defined in the space the artifact actually lives in; we did not do that here.
  \item \textbf{Matched-n does not neutralize generator quality.} In the instruction head-to-head (\S{}9.8) the released corpora are 20-70x larger than ours and were produced by different, mostly stronger generators, with human curation in at least one case. Matched-n sampling controls for size; it cannot control for the model that wrote the items, and PersonaHub's persona catalogue is orders of magnitude larger than any axis lattice we elicit.
\end{itemize}

\section{Conclusion}
At a finite budget, corpus construction is a covering problem, and treating it as one pays: the coverage functional is exactly submodular in its Monte-Carlo form, greedy selection is near-optimal in theory and measurement, and the resulting suites blanket the reachable space instead of decorating its boundary. Packing --- the right shape for an unbounded stream --- is measurably the wrong one here, sacrificing 2--4$\times$ covered fraction to win a min-gap metric no budgeted product requirement asks for. The sibling stack's scaffolding (elicited axes, gates, ledger, recursive refinement) transfers whole; what does not transfer is innocence about denominators: a method that grows the space it is covering must say so when it reports how much of that space it covered.

\bigskip\hrule\bigskip

\emph{All numbers in this paper are produced by \texttt{simulate\_coverage.py}, \texttt{verify\_theory.py}, and \texttt{live\_pilot.py} in this directory and recorded in \texttt{figures/summary\_sim.json}, \texttt{figures/summary\_theory.json}, and \texttt{figures/summary\_live.json}. Figures: \texttt{make\_figures.py}. Reproduction: \texttt{README.md}.}

\end{document}
